\subsection{Spectrum of an element in an algebra}

Let A be an algebra over K. Recall (A, III, p.~4) that one defines on the vector space \(\widetilde{\mathbf{A}} = \mathbf{K} \times \mathbf{A}\) an algebra structure such that:

\[
(\lambda , a) (\mu , b) = (\lambda \mu , \lambda b + \mu a + a b).
\]

Let \(e = (1, 0)\) . Then \((\tilde{\mathrm{A}}, e)\) is the unital algebra obtained from A by adjoining a unit element. The algebra A is identified with the two-sided ideal \(\{0\} \times A\) of \(\tilde{A}\) ; the algebra A is commutative if and only if \(\tilde{A}\) is.

If A' is a second algebra over K, let \((\tilde{\mathrm{A}}', e')\) be the unital algebra obtained from A' by adjoining an identity element, and let \(\varphi\) be a morphism from A to A'; there exists one and only one unital morphism from \((\tilde{\mathrm{A}}, e)\) into \((\tilde{\mathrm{A}}', e')\) which extends \(\varphi\) .

Let A be an algebra over K and \(x \in A\) . We call the spectrum of x relative to A the spectrum of x relative to the unital algebra \(\widetilde{A}\) deduced from A by adjoining an identity element. This set will be denoted \(\mathrm{Sp}_{\mathrm{A}}^{\prime}(x)\) , or \(\mathrm{Sp}^{\prime}(x)\) if no confusion arises. We have \(0 \in \mathrm{Sp}_{\mathrm{A}}^{\prime}(x)\) for all \(x \in A\) .

If \(\varphi\) is a morphism from A into an algebra B, we have \(\mathrm{Sp}_{\mathrm{B}}^{\prime}(\varphi(x)) \subset \mathrm{Sp}_{\mathrm{A}}^{\prime}(x)\) .

Remarks. --- 1) Let (A,1) be a unital algebra. If \(x \in A\) , one has

\[
\operatorname{Sp} _ {\mathrm{A}} ^ {\prime} (x) = \operatorname{Sp} _ {\mathrm{A}} (x) \cup \{0 \}.
\]

We verify in fact that \((e-1) \cdot A = A \cdot (e-1) = 0\) , so that \(\widetilde{A}\) is the unital algebra product of A and \(K(e-1)\) . Our assertion therefore follows from Remark 8 of I, p.~3.

\begin{enumerate}
\def\labelenumi{\arabic{enumi})}
\setcounter{enumi}{1}
\tightlist
\item
  It follows from Remark 1 that, if B is an algebra over K and if \(x \in B\) , we have:
\[
\mathrm{Sp} _ {\mathrm{B}} ^ {\prime} (x) = \mathrm{Sp} _ {\widetilde {\mathrm{B}}} (x) = \mathrm{Sp} _ {\widetilde {\mathrm{B}}} (x) \cup \{0 \} = \mathrm{Sp} _ {\widetilde {\mathrm{B}}} ^ {\prime} (x).
\]

\item
  If x belongs to the radical of A (A, VIII, p.~430, def. 3), then \(\mathrm{Sp}_{\mathrm{A}}^{\prime}(x)=\{0\}\). This follows from Remark 7 of I, p.~3.
\end{enumerate}

PROPOSITION 1. --- Let A be an algebra and \(x, y \in \mathrm{A}\) . We have \(\mathrm{Sp}'(xy) = \mathrm{Sp}'(yx)\) .

Passing to \(\widetilde{\mathbf{A}}\) we reduce to the case where A has a unit element \(e\) . It then suffices to prove that, if \(\lambda \neq 0\) is such that \(xy - \lambda e\) admits an inverse \(u\) , then \(yx - \lambda e\) is invertible. Set \(z = yux - e\) . Since \(xyu = \lambda u + e\) , one has

\[
\begin{array}{r l} (y x - \lambda e) z & = y (x y u) x - y x - \lambda y u x + \lambda e \\ & = y (\lambda u + e) x - y x - \lambda y u x + \lambda e = \lambda e \end{array}
\]

and likewise \(z(yx - \lambda e) = \lambda e\) . Since \(\lambda \neq 0\) , we see that \(yx - \lambda e\) is invertible.

If A is a unital algebra and if \(x, y \in A\) , the preceding proposition implies that \(\operatorname{Sp}(xy) \cup \{0\} = \operatorname{Sp}(yx) \cup \{0\}\) , but one may have \(\operatorname{Sp}(xy) \neq \operatorname{Sp}(yx)\) (cf.~I, p.~153, exerc. 3).

